\section{Actual finite tame type on the odd-weight branch}
\label{sec:cubic-actual-finite-type}

\begin{theorem}[AC-G78-type: qualitative finite-type occurrence]
\label{thm:G78-type}
Retain the inherited actual \(269\)-coefficient Type II hypotheses
of AC-G78, with \(e\) odd, the fixed cubic character \(\chi^j\),
and the host \(F/F^+\) of AC-G77.
In particular retain the original residual separation, ordinary
local conditions, and old modularity inputs used to construct
the endpoint \(V\).
Let \(q\) and the actual ordinary Steinberg occurrence be supplied
by AC-G78.  Choose a place \(v_0\) of \(F^+\) over \(q\), and the
chosen place of \(F\) above it; both have residue degree one.

There is an actual characteristic-zero automorphic representation
\(\sigma\) of the same definite unitary group, of the same symmetric
classical weight
\[
                         \mu(e)=(e-1,0,1-e),
\]
ordinary at the coefficient prime, and congruent to the original
cubic residual Hecke system away from the fixed bad places and \(q\),
such that its \(v_0\)-component contains the finite representation
\(\tau_\theta\) of AC-G72-character.
Here \(\theta\) is any fixed character of order \(269\) of
\(\mathbf F_{q^2}^{\times}\), trivial on
\(\mathbf F_q^\times\).  Such characters exist because
\(q\equiv-1\pmod{269}\).
Explicitly, let \(C_\theta^{(2)}\) denote the cuspidal
representation of \(\mathrm{GL}_2(\mathbf F_q)\) attached to
\(\theta\), let \(P_{2,1}\) be the standard parabolic subgroup
of \(\Gamma_q=\mathrm{GL}_3(\mathbf F_q)\), and put
\[
 \tau_\theta=
 \operatorname{Ind}_{P_{2,1}}^{\Gamma_q}
       \operatorname{Inf}_{\mathrm{GL}_2(\mathbf F_q)
                      \times\mathrm{GL}_1(\mathbf F_q)}^{P_{2,1}}
                    \bigl(C_\theta^{(2)}\otimes1\bigr).
 \tag{G78T.0}
\]
AC-G72-character proves that this is irreducible of dimension
\(q^3-1\).  All multiplicities below concern this
\(\mathrm{GL}_3(\mathbf F_q)\)-representation, not its
two-dimensional-Levi cuspidal factor.

Consequently the local conclusions of AC-G72-local apply:
the Weil--Deligne parameter at \(v_0\) is an irreducible tame
two-dimensional block, with the prescribed conjugate characters
\(\theta,\theta^q\), plus an unramified one-dimensional block,
up to the allowed unramified twists; its monodromy is zero.
The normalized global Galois representation is irreducible of
dimension three, by AC-G72-stable.
The ordinary condition and the symmetric weight are unchanged.

This is a qualitative occurrence theorem.  It gives no bound on
\(q\), does not descend to the original tame level or impose all
the original away-place local types, and does not prove the
per-curve Hecke-ideal containment in the modified cubic ring.
\end{theorem}

\begin{proof}
Write
\[
 K_q=\mathrm{GL}_3(\mathbf Z_q),\qquad
 K_q^1=\ker(K_q\longrightarrow\mathrm{GL}_3(\mathbf F_q)),
 \qquad \Gamma_q=K_q/K_q^1.
\]
The identification is the actual split local identification in
Thorne's group.  Fix the sufficiently small level away from
\(v_0\) used in AC-G78, and replace the hyperspecial level at
\(v_0\) by \(K_q^1\).
Choose a finite coefficient extension containing the required
finite-group character values and all the finitely many
characteristic-zero eigensystems that occur below.
Let \(\mathcal O\) be its ring of integers.
This auxiliary coefficient extension is used only to discuss
the automorphic lattice and its characters; it does not change
the retained coefficient-ring hypotheses of the original
deformation theorem.

\emph{The actual projective lattice.}
Use the algebraic weight lattice of weight \(\mu(e)\), and the
finite coefficient-prime level in the classical ordinary space
\(S_{\mu(e)}\).
The full definite-form lattice at \(K_q^1\) is a free
\(\mathcal O[\Gamma_q]\)-module.
Indeed, its finite double-coset set is a free \(\Gamma_q\)-set:
a stabilizer would give a rational element lying in the compact
level everywhere, hence a torsion element, and the level away
from \(v_0\) is sufficiently small.
Choosing a representative of each orbit identifies the associated
weight-lattice sections over that orbit with a regular group-ring
module tensored with the finite free weight lattice.
This is the literal definite-form argument of AC-G72, not an
assumption of exactness of taking \(269\)-group invariants.

The normalized ordinary operator at \(269\) is integral on this
lattice and commutes with \(\Gamma_q\).
Its ordinary projector is therefore an integral commuting
direct-summand projector.
Let \(\mathfrak m\) be the residual Hecke system of AC-G78,
using the unramified Hecke operators away from all the fixed bad
places and \(q\).
Their image on the finite ordinary lattice is a finite commutative
\(\mathcal O\)-algebra.  Since \(\mathcal O\) is complete, its
\(\mathfrak m\)-factor is cut out by an integral central idempotent.
This idempotent also commutes with \(\Gamma_q\).
Thus the actual fixed-weight, ordinary, residual-localized lattice
\(P\) is a finite projective \(\mathcal O[\Gamma_q]\)-module.

The Steinberg representation furnished by AC-G78 contributes to
this \(P\) after extension of scalars.
Its ordinary weight is the same \(\mu(e)\), its away Hecke system
reduces to \(\mathfrak m\), and an unramified twist of smooth
Steinberg has nonzero \(K_q^1\)-invariants.
Those invariants are the finite Steinberg representation
\(\mathrm{St}_{\Gamma_q}\).
No change of coefficient-prime refinement or of classical weight
has been made.

\emph{Genericity holds on the entire residual factor.}
AC-G78 supplies a good split witness place \(w_{\rm gen}\) for
the literal residual genericity condition of Thorne's
Theorem~5.2.
Every characteristic-zero automorphic constituent of \(P\)
has the same residual unramified Hecke polynomial at this place,
because \(w_{\rm gen}\) was kept outside the excluded set.
In particular every such constituent satisfies that same
residual genericity condition.

Here is the exact global reason this controls the component
at \(v_0\).  The paragraph following Thorne's Corollary~4.3,
printed p.~21, states that a nongeneric split local component
forces two irreducible subquotients of the global Galois
representation to differ by a cyclotomic twist.
After choosing stable lattices, these two constituents give
a pair of residual Frobenius eigenvalues whose ratio is the
cyclotomic value at \(w_{\rm gen}\).
This contradicts the fixed generic residual polynomial there.
Thus every split local component of every characteristic-zero
constituent of \(P\), in particular every \(v_0\)-component,
is generic.
No assertion that genericity only holds at the witness place,
and no Ramanujan estimate, is being substituted for this source
statement.

\emph{There is a strictly positive finite multiplicity difference.}
After extension of scalars, decompose the finite-dimensional
space \(P_E\) into its automorphic constituents.
As a \(\Gamma_q\)-representation it is a finite direct sum of
spaces \(\pi_q^{K_q^1}\), each with a nonnegative integer
multiplicity recording the other local factors and the chosen
ordinary subspace.
Old-vector multiplicities cause no problem: only these
nonnegative dimensions enter the argument.

For a generic spherical irreducible representation \(\pi_q\)
of \(\mathrm{GL}_3(\mathbf Q_q)\), its unramified principal-series
induction is itself irreducible, and
\[
 \pi_q^{K_q^1}\simeq
       \operatorname{Ind}_{B(\mathbf F_q)}^{\Gamma_q}1.
 \tag{G78T.1}
\]
One precise justification is Zelevinsky, Theorem~9.7(b),
printed p.~199, with its proof in Sections~9.8--9.10:
a generic irreducible is an irreducible induction of unlinked
essentially square-integrable segments.
For a spherical representation the cuspidal support consists
of unramified characters of \(\mathrm{GL}_1\).
Compact Iwasawa decomposition identifies the maximal-compact
invariants of that irreducible induction with the tensor product
of those of the segments.
An unramified Steinberg segment of length greater than one has
no such invariants.  Hence every segment has length one,
which proves the principal-series assertion.
Formula (G78T.1) then follows from the same compact Iwasawa
decomposition and reduction modulo \(q\).

Both the trivial representation and the finite Steinberg
representation occur with multiplicity one in (G78T.1).
For a nonspherical \(\pi_q\), the multiplicity of the trivial
\(\Gamma_q\)-representation in \(\pi_q^{K_q^1}\) is
\(\dim\pi_q^{K_q}=0\), whereas the finite Steinberg multiplicity
is nonnegative.
The actual smooth Steinberg constituent from AC-G78 contributes
one finite Steinberg and no trivial representation.
Summing these statements over the actual automorphic
decomposition gives the strict inequality
\[
 \bigl[P_E:\mathrm{St}_{\Gamma_q}\bigr]
                    -[P_E:1] >0.                   \tag{G78T.2}
\]
This proves positivity on the entire residual-localized space;
it does not infer positivity from one constituent while ignoring
possible negative contributions.

\emph{Projectivity converts the difference into an actual type.}
Since \(269\mid q+1\), fix \(\theta\) as in the statement.
AC-G72-character proves the Brauer-character identity
\[
             [\overline{\tau_\theta}]
                  =[\overline{\mathrm{St}_{\Gamma_q}}]-[1].
 \tag{G78T.3}
\]
The character of a projective \(\mathcal O[\Gamma_q]\)-lattice
vanishes on \(269\)-singular elements.
Taking the ordinary character inner product with \(P_E\),
(G78T.3), which is an identity on the \(269\)-regular elements,
therefore gives
\[
 [P_E:\tau_\theta]
       =\bigl[P_E:\mathrm{St}_{\Gamma_q}\bigr]-[P_E:1].
 \tag{G78T.4}
\]
Equivalently this is the projective Brauer reciprocity calculation
proved in AC-G72.
By (G78T.2) the right side is positive.
Consequently at least one actual automorphic constituent of
\(P_E\) contains \(\tau_\theta\) at \(v_0\).
This proves actual characteristic-zero occurrence, not only
a virtual-character or residual-lattice identity.

\emph{Local and global conclusions.}
The exact depth-zero type identification of AC-G72-local now
applies to this constituent.
It gives the stated irreducible tame two-dimensional block plus
unramified line and zero monodromy.
The source-normalized global representation is
\(r_{\sigma,\iota}\); put \(r=\epsilon\,r_{\sigma,\iota}\).
This removes the common rank-three Tate factor in (G78.1).
It leaves the local finite type unchanged, since \(\epsilon\)
is unramified at \(q\), and gives the unitary conjugate-self-dual
normalization used in AC-G72-stable.

For clarity the decisive residual placement can be checked
directly.  At \(q\) the residual roots of \(r\) are
\(\{-1,1,-1\}\).
The irreducible tame two-dimensional local block has Frobenius
eigenvalues \(a,-a\).
Since its finite inertia has order \(269\), it becomes trivial
in the residual semisimplification.
The residual multiset therefore forces
\(\overline a=\pm1\) and makes the remaining local line have
Frobenius value \(-1\).

Globally the residual characters of \(r\) are
\(\chi^j\overline\gamma,1,\chi^j\overline\gamma^{-1}\).
They are distinct, and only the middle character is individually
conjugate self-dual: \(\chi\) is anticyclotomic of order three,
\(\gamma\) is rational, and the retained residual inertia
condition gives \(\overline\gamma^2\ne1\).
These facts survive the host of AC-G77.
If the semisimple \(r\) had a two-dimensional constituent plus
a line, conjugate duality would fix that unique global line.
Its reduction would consequently be the middle residual
character \(1\).
But locally it must be the sole one-dimensional summand of the
prescribed type, whose residual Frobenius is \(-1\), a
contradiction.
A sum of three global lines cannot have an irreducible
two-dimensional local Weil--Deligne summand at all.
These exhaust the reducible dimension-three partitions,
including multiplicities.
Thus \(r\) is irreducible.
This is exactly the scope of AC-G72-stable, with the normalization
and residual placement made explicit.

All constructions took place in the fixed ordinary weight
and the same residual Hecke factor.
Only the local level at \(q\) was enlarged in obtaining the
finite type.
The proof supplies neither an effective bound for that prime
nor the additional original local-type and deformation-curve
identifications stated as remaining gates.
\end{proof}

\begin{remark}[Scope of later finite covers]
The controlled finite-cover results of AC-G73 can be applied
when their stated hypotheses are imposed, to remove the new
finite inertia while retaining irreducibility and the cubic
character.
That is a separate application; it does not supply a missing
bound for \(q\), control every other away-place local type,
or establish the original per-curve Hecke-ideal containment.
\end{remark}
